\section{Validity and Invalidity Across the Seven Sacraments}

An invalid sacramental attempt does not produce the sacrament. An illicit sacrament is truly effected but celebrated contrary to law. An abuse is a violation of liturgical or ecclesiastical discipline; some abuses are grave, but only a defect that touches what validity requires makes the act invalid. “Doubtful” describes the evidence available to a judge, not a third ontological outcome between valid and invalid. These categories are not evasions. They are how the Church protects sacramental certainty without calling every sin annihilation.

\subsection{The causal grammar}

Christ acts through an ecclesial sign. According to the sacrament, validity requires the divinely instituted sensible reality called matter or quasi-matter, the determinative sacramental words or form, a minister capable of acting, intention at least to do what the Church does, and a capable recipient. Matrimony and Penance require further acts intrinsic to their signs. The Church determines and approves rites around this substance and can authoritatively identify what its integrity demands; she cannot abolish what Christ instituted.

The minister need not understand every theological implication, possess grace, or personally believe every dogma. If personal sin invalidated a sacrament, certainty would collapse into Donatism. Neither is intention a magical safety net for arbitrary wording. The DDF declaration \emph{Gestis verbisque} (2024) insists that matter and form belong together in the Church's act and that grave unauthorized alteration of an essential formula can prevent the sacrament. Fidelity to approved words is therefore pastoral charity, not rubricist anxiety.

Approved rites carry the Church's public judgment and a presumption of validity when enacted by a capable minister. The 1962 rites and the revised Roman rites promulgated by Paul VI and his successors are not rival sacramental religions. Their different formulas and ceremonies are valid ecclesial determinations of the same seven sacraments. A person claiming that the revised rites of Confirmation, Anointing, or Orders are invalid bears more than a historical-aesthetic burden: the claim opposes the competent Apostolic constitutions that identified their essentials and the Church's universal sacramental reception.

\subsection{A seven-sacrament matrix}

\begin{longtable}{>{\bfseries\raggedright\arraybackslash}p{0.13\linewidth}>{\raggedright\arraybackslash}p{0.27\linewidth}>{\raggedright\arraybackslash}p{0.29\linewidth}>{\raggedright\arraybackslash}p{0.23\linewidth}}
\toprule
\textbf{Sacrament} & \textbf{What validity centrally requires} & \textbf{Typical invalidating defect} & \textbf{1962 / revised-rite judgment}\\
\midrule
\endhead
Baptism & True natural water applied by washing and the singular Trinitarian form---the minister baptizes this person in the name of Father, Son, and Holy Spirit---with intention to do what the Church does. In necessity any person can minister. & No true water; omission or substitution of a divine Person; a formula whose subject makes the community rather than Christ the minister, such as “We baptize you”; deliberate intention to perform only a social naming. & The older and current Roman rites are valid. Their exorcisms, anointings, interrogations, and ceremonies differ without replacing water and Trinitarian invocation. The CDF in 2020 declared the “we baptize” formula invalid.\\
Confirmation & A baptized, unconfirmed subject; validly consecrated chrism, anointing with laying on of the hand, the approved form, and a bishop or a priest who has the faculty required to confirm validly in the Latin Church. & No chrism or sacramental anointing; a substantially defective form; an unordained minister; a Latin priest lacking the needed faculty outside a provision of law. & Paul VI's \emph{Divinae consortium naturae} (1971) defined the revised rite and form. The traditional and revised forms validly seal with the gift of the Holy Spirit. Difference from the older wording is not defect.\\
Eucharist & Validly ordained bishop or priest; wheat bread and natural grape wine; the approved Eucharistic Prayer with the institutional words and ecclesial intention to consecrate and offer as the Church does. & Unordained “presider”; rice or other non-wheat substitute; liquid not grape wine; simulation or positive exclusion of consecration; substantial corruption of the consecratory form. & A Mass according to either Missal is capable of validity. Language, orientation, chant, altar arrangement, offertory wording, Eucharistic-prayer choice, reception posture, or a minister's sin does not of itself invalidate. Low-gluten hosts retaining wheat are valid; wholly gluten-free matter is not. Approved mustum is grape matter.\\
Penance & A validly ordained priest with faculty to absolve validly; a penitent capable of the sacrament who manifests sins and has at least imperfect contrition with a purpose of amendment; the priest's sacramental absolution. & Unordained confessor; priest without the faculty required by law, unless the law itself supplies it or danger-of-death law applies; no absolution; no contrition or intention to renounce mortal sin. & Older and current rites are valid. Absolution is judicial and medicinal, not merely declaratory counseling. Acts of the penitent are integral; exact individual culpability is not reducible to the ritual's outward appearance.\\
Anointing of the Sick & A seriously ill or frail baptized person with at least habitual implicit intention; validly ordained bishop or priest; blessed oil of olives or, under current law, another approved plant oil; anointing and the approved sacramental prayer. & Deacon or lay “anointer”; no oil/anointing; essential form omitted or substantially replaced; subject certainly dead or not capable of receiving. & Paul VI's \emph{Sacram unctionem infirmorum} (1972) revised matter, form, and administration within ecclesial competence. Both rites are valid. The sacrament is for the sick, not only the final minutes of life.\\
Holy Orders & A baptized male subject; validly ordained consecrating bishop; imposition of hands and the essential consecratory prayer appropriate to diaconate, presbyterate, or episcopate; intention to ordain as the Church ordains. & Unconsecrated ordainer; no laying on of hands; essential prayer substantially defective; a woman as recipient; positive intention to enact something other than sacramental Order. & Paul VI's \emph{Pontificalis Romani} (1968) designated the revised matter and essential forms. The revised rites validly ordain. Lack of papal mandate for an episcopal consecration makes it gravely illicit and penal, not by that fact invalid.\\
Matrimony & In the Latin Church, a capable baptized man and woman give irrevocable present consent to each other, free of invalidating impediment; Catholics ordinarily must observe canonical form before the competent or delegated assistant and two witnesses unless law or dispensation supplies another route. & No true consent; simulation or exclusion of marriage itself or an essential property; invalidating impediment without dispensation; defect of canonical form where the parties are bound; incapable parties. & The spouses are the ministers in the Latin theological account; the cleric receives consent in the Church's name. Old or new nuptial ritual does not by preference determine validity. Canonical form, delegation, impediments, and consent do. Eastern discipline has its own governing law and is not collapsed into this table.\\
\bottomrule
\end{longtable}

\subsection{Why revised forms do not create revised sacraments}

The Church has repeatedly specified essential signs while revising the surrounding rite. Pius XII did so for Orders in \emph{Sacramentum ordinis}; Paul VI did so for the revised rites of Orders, Confirmation, and Anointing. A form can express the same sacramental determination without repeating every phrase of a former form. Validity is not measured by verbal length or by resemblance to a preferred medieval book.

This authority is not unlimited. The Apostolic See could identify the laying on of hands and consecratory prayer within the received sacrament of Orders; it could not authorize ordination without a bishop or of a subject incapable by divine constitution. It could regulate oil and anointings; it could not make lay anointing a sacrament. It could approve Eucharistic prayers; it could not substitute a communal meal over rice and water for the Eucharist. The strongest defense of revised rites therefore includes, rather than dissolves, sacramental boundaries.

\subsection{The principal rite-by-rite disputes}

\textbf{Baptism.} The revised rite did not replace water or the command of Matthew~28:19. It reorders exorcisms, anointings, interrogations, and community roles around the same washing and singular Trinitarian act. The 2020 CDF judgment on “We baptize you” shows that postconciliar authority does not treat creative paraphrase as harmless: the communal subject obscures Christ acting through the individual minister and invalidates the attempt. By contrast, an emergency Baptism with true water and the exact Trinitarian form can be valid without the surrounding ceremonies; they should later be supplied as law directs, but they do not retroactively cause the sacrament.

\textbf{Confirmation.} The older Latin form and Paul VI's \latin{Accipe signaculum doni Spiritus Sancti} do not need identical word counts. The apostolic constitution identifies anointing with chrism on the forehead, performed by laying on the hand, and the approved form as the sign that seals with the Spirit's Gift. Eastern Catholic forms already demonstrate that one divine sacrament is not imprisoned in one medieval Latin sentence. The substantial question in a concrete Latin case is not “old or new?” but whether the minister, chrism, anointing, form, intent, subject, and priestly faculty converge.

\textbf{Eucharist.} The offertory prayers do not consecrate; replacing the medieval offertory therefore cannot make the new Eucharist invalid. Nor do Latin, eastward orientation, a silent Canon, or the number of signs of the cross belong to validity. The approved Eucharistic Prayers contain the institutional words within a consecratory prayer offered by a valid priest over wheat bread and grape wine. Eucharistic Prayer II's brevity and disputed ancient ancestry are legitimate comparative questions, not invalidating defects.

The former approved English rendering “for you and for all” was theologically and translationally contested and was corrected to “for you and for many” in 2011. It did not make decades of approved English Masses invalid: the Holy See approved and interpreted the formula within the complete institutional narrative and Catholic intention. The distinction between Christ's universally sufficient redemption and the fruit received by many remains doctrinally important; a deficient translation can require correction without becoming a lay tribunal's declaration that no consecration occurred.

\textbf{Penance.} Both rites culminate in the priest's judicial sacramental act, whose essential Latin determination includes \latin{Ego te absolvo}. The newer form surrounds it with an explicitly Trinitarian prayer. A communal reconciliation service without individual confession and absolution is not the sacrament merely because it is held in a church. General absolution can be valid only under the law's grave conditions and with the recipient's required disposition, including intention to confess grave sins individually in due time. A valid priest without the needed faculty ordinarily cannot absolve validly, however traditional his stole or formula.

\textbf{Anointing.} Paul VI reduced and reordered anointings, broadened the immediate pastoral use beyond an almost exclusively terminal setting, and promulgated a form asking that the Lord in mercy help the sick person with the grace of the Holy Spirit, save, and raise him. Those changes remain within the sacrament witnessed in James~5:14--15. A deacon's or layperson's compassionate anointing is not the sacrament; an old priestly rite and the revised priestly rite are.

\textbf{Orders.} Objections to the revised ordinals often isolate one omitted medieval phrase while ignoring the consecratory prayer Paul VI designated as the form and the laying on of hands he designated as matter. The episcopal prayer asks the Spirit who governs and guides, and the presbyteral prayer asks renewal of the Spirit of holiness for priestly office within an ordaining Church that unmistakably intends Catholic Order. Apostolic succession, a valid consecrating bishop, a capable male subject, imposition of hands, form, and intention remain. A woman cannot be ordained; an actor cannot self-ordain; and a papally unmandated but validly consecrated bishop may sacramentally ordain while committing a gravely unlawful act.

\textbf{Matrimony.} The nuptial Mass can surround consent with different readings, blessings, and ritual placements, but the spouses' covenantal consent effects the bond in the Latin account. Changing the nuptial blessing does not change the ministers. The hard validity cases concern consent, capacity, impediments, and canonical form: simulation, exclusion of permanence or fidelity, coercion, a prior bond, or an unauthorized witness where canonical form binds. An exquisitely celebrated old rite cannot cure those defects, and a simply celebrated approved new rite does not create them.

\subsection{Ministerial irregularity does not have one sacramental answer}

“His orders are valid” never means “all his sacraments are valid.” Faculties, jurisdiction, and canonical form affect different sacraments differently.

\begin{fourstable}{Situation}{Validity in the ordinary case}{Liceity / communion}{Controlling qualification}
A lay person baptizes with water, Trinitarian form, and ecclesial intention & Valid, even though outside necessity the action can be gravely illicit. & The minister does not become clergy and cannot minister other sacraments. & Baptism's emergency minister is exceptionally broad.\\
A valid priest celebrates Mass without authorization & If matter, form, and intention are present, the Eucharist is valid. & Gravely illicit where he lacks mission or violates liturgical law; attendance raises separate ecclesial and moral questions. & Eucharistic validity does not supply incardination, office, or communion.\\
A valid priest hears confession without faculty & Ordinarily invalid. & Also illicit. & Canons 966 and 976, supplied faculty under canon 144 where truly applicable, and other provisions must be assessed from complete facts.\\
A Latin priest confirms without a faculty & Invalid in the ordinary case. & The law itself grants faculties in specified situations; an undocumented assumption is unsafe. & Priesthood alone is not the complete Latin validity condition for Confirmation.\\
A valid priest anoints without a lawful assignment & Sacramentally valid if the other essentials are present. & May be illicit and pastorally disordered. & Only bishops and priests can validly anoint; a deacon cannot.\\
A valid bishop ordains without papal mandate & Capable of being valid if subject, rite, and intention are sound. & Gravely illicit and subject to the penalties in current law; it can constitute a schismatic act. & Mandate and canonical mission are not the matter or form of Orders.\\
A Catholic attempts marriage before an unauthorized cleric & Ordinarily invalid for defect of canonical form if no delegation, dispensation, supplied faculty, or extraordinary form applies. & The parties' subjective good faith does not itself repair form. & Matrimonial validity is fact-intensive; a tribunal or competent authority judges concrete cases.\\
\end{fourstable}

\subsection{The SSPX after the acts of July 2026}

Historical descriptions of the Society of Saint Pius X as “canonically irregular but not in schism” are no longer a safe current endpoint. After the Society proceeded with episcopal consecrations despite the Holy See's warning and Leo XIV's prohibition, the DDF explanatory note of 2 July 2026 applied the current canonical consequences. It states that the consecrating and consecrated bishops incurred excommunication and that SSPX ministers are now in schism; their sacramental celebrations are illicit. The note further states that Penance and marriages celebrated with SSPX ministry are invalid under the present situation. The faithful are asked to abstain from SSPX sacramental and liturgical events; merely attending does not mechanically establish the canonical crime of schism, which requires formal adherence.

The same act does \emph{not} declare every SSPX priest unordained or every SSPX Eucharist invalid. Where a validly ordained priest uses valid matter, form, and intention, Eucharistic validity remains a sacramental question distinct from schism and liceity. Earlier special faculties for confession and arrangements for marriage cannot simply be quoted as though the July 2026 act had not changed the controlling facts. Current law must be read from its current endpoint.

\subsection{Independent, sedevacantist, and separated bodies}

No global verdict can responsibly validate every sacrament claimed by an independent network. Apostolic lineage must be documented; the ordaining rite and intention must be known; each minister and act must be evaluated; Penance faculties and Catholic matrimonial form present separate barriers. A body may possess valid episcopal orders while lacking canonical mission. Another may have broken or doubtful lineage. Self-certification is not competent ecclesiastical judgment.

Sedevacantism also adds a doctrinal rupture: it treats the postconciliar popes and often the revised sacramental rites as invalid or doubtful. That assertion is not made true by the antiquity of the liturgy used. A traditional form can be used to enact a non-Catholic ecclesiology. Conversely, a Catholic celebrant's use of the postconciliar form does not make his Orders doubtful.

\subsection{Recipient, fruit, and sacrilege}

The sacraments work \latin{ex opere operato}: Christ is faithful when the sign is validly enacted. They do not sanctify a resistant recipient mechanically. Baptism and Orders validly confer character even when received illicitly; sacramental Communion by one conscious of grave sin can be objectively valid and personally sacrilegious; absolution cannot justify a person who withholds contrition. Validity concerns whether Christ's sacramental act occurs. Fruitfulness concerns how grace is received. Culpability concerns knowledge, freedom, and consent.

If a validity defect is credibly suspected, the answer is not Internet diagnosis or casual repetition. Records, exact formulas, matter, ordination lineage, faculty, canonical form, and competent judgment must be gathered. Conditional administration has a disciplined place where a positive prudent doubt cannot otherwise be resolved; it is not a ritual expression of generalized mistrust.

\canonwarning{This chapter states universal Latin principles as of 13 July 2026. It cannot determine an individual's Baptism, Orders, absolution, marriage, excommunication, or Sunday obligation. Such cases require the complete facts, current particular and proper law, and the competent ordinary, tribunal, dicastery, or qualified canonist.}
